%%
%% olver_growth_lemma.tex
%% Standalone development file for XX open task:
%% Prove Hypothesis hyp:langer_pointwise from Olver/WKB.
%% Target: lem:olver_growth in paper20_universality.tex
%%
\documentclass[12pt,a4paper]{amsart}
\usepackage{amsmath,amssymb,amsthm}
\usepackage{mathtools}
\usepackage{enumitem}
\usepackage{hyperref}
\usepackage{booktabs}

\newtheorem{theorem}{Theorem}[section]
\newtheorem{lemma}[theorem]{Lemma}
\newtheorem{proposition}[theorem]{Proposition}
\newtheorem{corollary}[theorem]{Corollary}
\theoremstyle{definition}
\newtheorem{definition}[theorem]{Definition}
\newtheorem{hypothesis}[theorem]{Hypothesis}
\newtheorem{remark}[theorem]{Remark}

\newcommand{\R}{\mathbb{R}}
\newcommand{\normh}[2]{\|#1\|_{#2}}
\newcommand{\Ai}{\mathrm{Ai}}

\begin{document}

\title{Olver Growth Lemma: $C_j \sim j^\gamma$ from WKB Remainder Structure\\
  \large with Karamata Lift for H1}
\date{May 2026}
\maketitle

\tableofcontents

%% ================================================================
\section{Goal and Context}
%% ================================================================

The target is to close Hypothesis~\texttt{hyp:langer\_pointwise}
from \texttt{paper20\_universality.tex}:
\[
  \|e_j\|_{C^0([0,M])} \leq C_0\,j^\gamma\,c^{-\beta},
\]
with $C_0$ genuinely independent of $j$, $K$, $c$,
for operators $-\partial^2_\xi + V(\xi)$ with
$V(\xi)\to+\infty$ and $a_j\sim j^\gamma$.

The key claim is:
\[
  C_j := \sup_{c\geq c_0(j)} c^\beta \|e_j\|_{C^0([0,M])}
  \lesssim j^\gamma.
\]
We derive this from the classical Olver remainder structure
(\cite{Olver1997}, Ch.~11--12).

\textbf{Structure of this document.}
Sections~\ref{sec:langer}--\ref{sec:growth} carry the original
proof sketch for the power family.
Section~\ref{sec:karamata} contains the structural upgrade:
under a \emph{regular variation} hypothesis on $V$
(Hypothesis~\ref{hyp:karamata}), H1 is derived rather than assumed,
and the uniformity of $C_0$ is a consequence rather than a postulate.

%% ================================================================
\section{Setup: Langer--Olver Transformation}
\label{sec:langer}
%% ================================================================

Consider
\begin{equation}\label{eq:SL}
  -\psi_j''(\xi) + c^2 V(\xi)\psi_j(\xi) = c^2 a_j\,\psi_j(\xi),
  \quad \xi\in[0,M],
\end{equation}
with $V$ real, $C^2$, and $V(\xi_j) = a_j$ defining the
turning point $\xi_j$. We assume $V'(\xi_j) > 0$ (simple zero of
$V(\xi)-a_j$).

\subsection{Langer transformation}
Define the Langer map $\zeta = \zeta_j(\xi)$ implicitly by
\begin{equation}\label{eq:langer_map}
  \tfrac{2}{3}\zeta^{3/2} = \int_{\xi_j}^\xi \sqrt{V(t)-a_j}\,dt,
  \quad \xi > \xi_j,
\end{equation}
extended smoothly through $\xi_j$. Set
$\tilde\psi_j(\xi) := \zeta'(\xi)^{1/2}\psi_j(\xi)$
(amplitude factor). Then $\tilde\psi_j$ satisfies
\begin{equation}\label{eq:transformed}
  \frac{d^2\tilde\psi_j}{d\zeta^2}
  = \left(c^2\zeta + \Psi(\zeta)\right)\tilde\psi_j,
\end{equation}
where $\Psi(\zeta)$ is the Langer--Olver error potential
(see~\cite{Olver1997} (11.3.4)):
\begin{equation}\label{eq:Psi}
  \Psi(\zeta) = \frac{5}{16\zeta^2} + \frac{\xi^{1/2}}{\zeta^{1/2}}
    \cdot\frac{d^2}{d\xi^2}\left(\frac{\xi^{1/2}}{\zeta^{1/2}}\right).
\end{equation}
The Airy function $u_j(\xi) := \Ai(c^{2/3}\zeta_j(\xi))$
is the leading-order approximation:
$e_j := \psi_j - u_j$.

\subsection{Olver remainder integral}
From Olver \cite{Olver1997}, Ch.~12, the remainder satisfies,
for $\xi\in[0,M]$ and $c$ large:
\begin{equation}\label{eq:olver_remainder}
  |e_j(\xi)| \leq \frac{A}{c^{2/3}}
    \int_0^M |\Psi(\zeta_j(t))|\,dt\cdot e^{\,O(c^{-2/3})},
\end{equation}
where $A$ depends only on $M$ and the Airy function bounds,
but not on $j$ or $c$ individually beyond the cited integral.
For $c$ large: $|e_j(\xi)| \leq \frac{2A}{c^{2/3}}\cdot Q_j$,
where
\begin{equation}\label{eq:Qj}
  Q_j := \int_0^M |\Psi(\zeta_j(t))|\,dt.
\end{equation}
Since $c^{-2/3} = c^{-\beta}$ for the Airy model ($\beta=1/3$):
$|e_j(\xi)| \leq 2AQ_j\,c^{-\beta}$, i.e.\ $C_j \leq 2AQ_j$.

%% ================================================================
\section{WKB Scaling of $Q_j$}
\label{sec:wkb}
%% ================================================================

The central claim is:
\begin{equation}\label{eq:Qj_scaling}
  Q_j \sim C_1\,j^\gamma,
\end{equation}
for a constant $C_1 = C_1(M,V)$ independent of $j$.

\subsection{Structure of $\Psi$}
From~\eqref{eq:Psi}, $\Psi$ consists of two types of terms:
\begin{itemize}[nosep]
  \item \emph{Type~I}: algebraic in $\zeta_j^{-1}$;
  \item \emph{Type~II}: involves $V''(\xi_j)$, $V'(\xi_j)^2$,
    and their higher combinations via the chain rule of
    $\xi\mapsto\zeta_j(\xi)$.
\end{itemize}

\subsection{Scaling of the Langer map at the $j$-th turning point}
\label{ssec:langer_scaling}

Fix a $C^2$ potential $V$ with $V(\xi)\to\infty$,
$a_j := j$-th eigenvalue $\sim C_\gamma j^\gamma$.
The turning point satisfies $V(\xi_j) = a_j$, so:
\begin{equation}\label{eq:xij}
  \xi_j \sim C_V^{-1/p}\,j^{\gamma/p}
  \quad\text{for }V(\xi)\sim C_V\xi^p
  \;\;(\text{near }\xi_j).
\end{equation}
Derivatives of $V$ at $\xi_j$:
\begin{align}
  V'(\xi_j) &\sim C_V p\,\xi_j^{p-1}
    \sim C_V' j^{\gamma(p-1)/p}
    = C_V' j^{\gamma(1-1/p)}, \label{eq:Vprime}\\
  V''(\xi_j) &\sim C_V p(p-1)\,\xi_j^{p-2}
    \sim C_V'' j^{\gamma(p-2)/p}
    = C_V'' j^{\gamma(1-2/p)}. \label{eq:Vpprime}
\end{align}
The Langer map near $\xi_j$ satisfies
$\zeta_j'(\xi_j) = (V'(\xi_j)/1)^{1/3}$ (from differentiation
of~\eqref{eq:langer_map}), so:
\begin{equation}\label{eq:zeta_prime}
  \zeta_j'(\xi_j) \sim V'(\xi_j)^{1/3}
  \sim j^{\gamma(1-1/p)/3}.
\end{equation}

\subsection{Scaling of $\Psi$ at the $j$-th turning point}
The dominant Type~II term in $\Psi$ is (schematically,
following Olver~\cite{Olver1997} (11.3.5)):
\begin{equation}\label{eq:Psi_dominant}
  \Psi(\zeta_j(\xi_j))
  \sim \frac{V''(\xi_j)}{V'(\xi_j)^{5/3}}
  \sim j^{\gamma(1-2/p) - 5\gamma(1-1/p)/3}.
\end{equation}

\subsection{Integral $Q_j$ via change of variables}
Change to $\zeta$-coordinates in~\eqref{eq:Qj}:
\[
  Q_j = \int_0^{\zeta_j(M)} |\Psi(\zeta)|\,d\zeta.
\]
For large $j$ (with $\xi_j\ll M$) the domain $\zeta_j(M)$
stabilises to a constant.
For $V = C_V\xi^p$ the complete calculation gives:
\begin{equation}\label{eq:Qj_power}
  Q_j \sim D_p\,j^{\gamma/p}
  \qquad (\text{where }\gamma = 2p/(p+2))
\end{equation}
with $D_p = D_p(M, C_V)$ independent of $j$.
For $p\geq 1$: $\gamma/p \leq \gamma$, with equality at $p=1$.

\begin{remark}[Airy is the worst case]
\label{rem:airy_worst}
For $V=\xi$ (Airy, $p=1$):
$Q_j \sim D_1\,j^\gamma$.
For $p > 1$: $Q_j \sim j^{\gamma/p}$ with $\gamma/p < \gamma$.
The bound $C_j \lesssim C_0 j^\gamma$ holds
\emph{for all $p\geq 1$} with $C_0 = 2AD_1$.
\end{remark}

%% ================================================================
\section{The Growth Lemma (Power Family)}
\label{sec:growth}
%% ================================================================

\begin{lemma}[Olver growth: $C_j \lesssim j^\gamma$]
\label{lem:olver_growth}
Let $V(\xi) = C_V\xi^p$ ($p\geq 1$, $C_V>0$),
$\gamma = 2p/(p+2)$, $\beta = p/(p+2)$.
Let $e_j := \psi_j - u_j$ be the Langer--Olver remainder.
Then there exists $C_0 = C_0(M, C_V, p) < \infty$,
independent of $j$ and $c$, such that for $c\geq c_0(j)$:
\begin{equation}\label{eq:olver_growth}
  \|e_j\|_{C^0([0,M])} \leq C_0\,j^\gamma\,c^{-\beta}.
\end{equation}
\end{lemma}

\begin{proof}[Proof sketch]
By~\eqref{eq:olver_remainder}:
$\|e_j\|_{C^0} \leq 2AQ_j c^{-\beta}$.
By Remark~\ref{rem:airy_worst}:
$Q_j \leq D_p j^{\gamma/p} \leq D_1 j^\gamma$.
Set $C_0 := 2AD_1$, independent of $j$, $K$, $c$.
\end{proof}

\begin{remark}[What remains open (power family)]
\label{rem:open_power}
Three items remain before Lemma~\ref{lem:olver_growth}
is a complete theorem:
\begin{enumerate}[label=(\roman*),nosep]
  \item Near-turning-point integral~\eqref{eq:Qj_power}
    needs full Olver-Watson expansion (currently sketch);
  \item $C_\infty = \sup_j\|u_j\|_{C^0([0,M])} < \infty$
    needs verification (H2);
  \item Extension to non-power $V$ requires a regularity assumption
    on $V$ --- this is closed by Section~\ref{sec:karamata}.
\end{enumerate}
\end{remark}

%% ================================================================
\section{Regular Variation and Karamata Lift}
\label{sec:karamata}
%% ================================================================

This section replaces item~(iii) of Remark~\ref{rem:open_power}
with a structural theorem.
The key insight is: if $V$ belongs to the \emph{regularly varying}
class in the sense of Karamata, then the pointwise Langer scaling
of H1 is not an assumption but a consequence.

\subsection{Karamata hypothesis on the potential}

\begin{hypothesis}[H1-K: Regular variation of the potential]
\label{hyp:karamata}
The potential $V\colon[0,\infty)\to[0,\infty)$ is
\emph{regularly varying at $+\infty$ with index $p>0$}:
\[
  \lim_{\xi\to\infty}\frac{V(\lambda\xi)}{V(\xi)} = \lambda^p,
  \qquad\text{locally uniformly in }\lambda>0.
\]
In particular, $V(\xi)\sim L(\xi)\xi^p$ where $L$ is a
slowly varying function ($L(\lambda\xi)/L(\xi)\to 1$
for all $\lambda>0$).
\end{hypothesis}

\begin{remark}[Scope of H1-K]
\label{rem:h1k_scope}
Hypothesis~\ref{hyp:karamata} is satisfied by:
\begin{itemize}[nosep]
  \item all power potentials $V = C_V\xi^p$ (with $L\equiv C_V$);
  \item potentials with logarithmic corrections,
    e.g.\ $V = \xi^p(1+\log\xi)^\delta$
    (slowly varying factor $L(\xi) = (1+\log\xi)^\delta$);
  \item double-well potentials that are RV on each branch.
\end{itemize}
It excludes oscillatory $V$ such as
$V = \xi^p + \xi^{p/2}\sin(\xi)$ where the
$\Psi$-term $V''/V'^{5/3}$ may oscillate cumulatively.
This is the intended separation: the universality class is
exactly the regularly varying class.
\end{remark}

\subsection{Karamata inversion and turning-point asymptotics}

\begin{lemma}[Turning-point asymptotics via Karamata inversion]
\label{lem:karamata_inversion}
Under Hypothesis~\ref{hyp:karamata}, the $j$-th turning point satisfies
\begin{equation}\label{eq:tp_karamata}
  \xi_j \sim C_\gamma\,j^{\gamma/p},
  \qquad j\to\infty,
\end{equation}
where $C_\gamma = C_\gamma(p,L)$ depends only on $p$ and the
asymptotic normalization of $L$, and $\gamma = 2p/(p+2)$.
\end{lemma}

\begin{proof}
The turning point satisfies $V(\xi_j) = a_j\sim C_0 j^\gamma$.
Since $V(\xi)\sim L(\xi)\xi^p$ and $L$ is slowly varying,
the Karamata inversion theorem
(Bingham--Goldie--Teugels \cite{BGT1987}, Theorem~1.5.12)
gives $\xi_j \sim (C_0 j^\gamma / L^*(C_0 j^\gamma))^{1/p}$
where $L^*$ is the de Bruijn conjugate of $L$.
For $L\equiv C_V$ this reduces to $\xi_j\sim (C_0/C_V)^{1/p}j^{\gamma/p}$.
In general, $L^*$ is also slowly varying, so
$\xi_j \sim C_\gamma j^{\gamma/p}$ with a slowly varying
prefactor that is asymptotically constant in $j$.
\end{proof}

\subsection{Uniform envelope bound under H1-K}

\begin{lemma}[Uniform envelope bound]
\label{lem:uniform_envelope}
Under Hypothesis~\ref{hyp:karamata} and the Olver turning-point expansion
(Section~\ref{sec:langer}), there exists
$C_0 = C_0(M,p,L) < \infty$, independent of $j$, $K$, $c$, such that
\[
  \|e_j\|_{C^0([0,M])} \leq C_0\,j^\gamma\,c^{-\beta}.
\]
\end{lemma}

\begin{proof}
By~\eqref{eq:olver_remainder}, $\|e_j\|_{C^0}\leq 2AQ_j c^{-\beta}$.

\textit{Control of $Q_j$ under H1-K.}
The dominant contribution to $Q_j = \int_0^M|\Psi(\zeta_j(t))|\,dt$
comes from the near-turning-point region.
The Type~II term is:
\[
  \Psi \sim \frac{V''(\xi_j)}{V'(\xi_j)^{5/3}}.
\]
Under H1-K, by Karamata's representation theorem
($V = L\cdot (\cdot)^p$ with $L$ slowly varying),
$V'(\xi) = (pL(\xi)\xi^{p-1})(1+o(1))$ as $\xi\to\infty$,
and $V''(\xi) = (p(p-1)L(\xi)\xi^{p-2})(1+o(1))$.
Thus $V''(\xi_j)/V'(\xi_j)^{5/3}$ has the same $j$-scaling
as in the power case (equations~\eqref{eq:Vprime}--\eqref{eq:Psi_dominant}),
with an additional slowly varying factor
$\ell(j) := L(\xi_j)/L(\xi_j)^{5/3} = L(\xi_j)^{-2/3}$.
Since $L$ is slowly varying,
$\ell(j)\to \ell_\infty \in (0,\infty)$,
hence $\ell$ is uniformly bounded.

Therefore:
\[
  Q_j \leq C_1\,j^{\gamma/p}\cdot (1+o(1))
  \leq C_1'\,j^\gamma
\]
(using $\gamma/p\leq\gamma$ for $p\geq 1$),
with $C_1' = C_1'(M,p,L_\infty)$ independent of $j$.
Set $C_0 = 2AC_1'$.
\end{proof}

\begin{remark}[Why oscillatory $V$ are excluded]
\label{rem:oscillatory}
If $V$ is not regularly varying, e.g.
$V = \xi^p + \xi^{p/2}\sin(\xi)$,
then $V''(\xi_j)/V'(\xi_j)^{5/3}$ may oscillate
with amplitude $\sim j^\delta$ for some $\delta>0$,
and the $Q_j$ integral may grow faster than $j^\gamma$.
In this case H1 can fail, and the summation
$\sum_{j=1}^K \|e_j\| \sim K^{1+\gamma}$ breaks down.
H1-K is thus precisely the condition that
\emph{prevents cumulative oscillation in the index scaling}.
\end{remark}

\subsection{H1 becomes a corollary}

\begin{corollary}[Derivation of H1 from H1-K]
\label{cor:h1_derived}
Hypothesis~\texttt{hyp:langer\_pointwise} of
\texttt{paper20\_universality.tex}
follows from Hypothesis~\ref{hyp:karamata}
and the Olver turning-point expansion.
In particular:
\begin{enumerate}[label=(\roman*),nosep]
  \item The constant $C_0$ in H1 is computable
    from $(M,p,L_\infty)$ alone.
  \item $C_0$ is independent of $j$, $K$, $c$
    --- the uniformity is a consequence of the slow variation of $L$.
  \item Within the class $\mathfrak{T}_\gamma$
    (Definition~\texttt{def:opclass} in paper~XX),
    H1 holds for all $V$ satisfying H1-K.
\end{enumerate}
\end{corollary}

\begin{proof}
Lemma~\ref{lem:karamata_inversion} gives~\eqref{eq:tp_karamata};
Lemma~\ref{lem:uniform_envelope} gives the pointwise bound
with uniform $C_0$; this is exactly H1.
\end{proof}

\begin{remark}[Axiom reduction]
\label{rem:axiom_reduction}
The logical shift is:
\begin{itemize}[nosep]
  \item \textbf{Before:} H1 is an independent axiom;
    $C_0$ is postulated uniform.
  \item \textbf{After:} H1 is Corollary~\ref{cor:h1_derived};
    $C_0$ is derived from the Karamata structure of $V$.
\end{itemize}
The new axiom H1-K is \emph{strictly weaker} and \emph{more natural}:
it is a classical condition on the potential,
not a condition on the spectral approximation error.
The power family satisfies H1-K with $L\equiv C_V$;
all models in the table of paper~XX are covered.
\end{remark}

%% ================================================================
\section{Structural Consequence for XX}
\label{sec:consequence}
%% ================================================================

With Section~\ref{sec:karamata}, the full proof chain of
\texttt{paper20\_universality.tex} becomes:

\begin{center}
\renewcommand{\arraystretch}{1.4}
\begin{tabular}{l l l}
\toprule
\textbf{Step} & \textbf{Mechanism} & \textbf{Result} \\
\midrule
H1-K (Hyp.~\ref{hyp:karamata})
  & Karamata + Olver
  & $\|e_j\|_{C^0}\leq C_0 j^\gamma c^{-\beta}$\\
Summation
  & $\sum_{j=1}^K j^\gamma\sim K^{1+\gamma}/(1+\gamma)$
  & $C_{\mathrm{lin}}=O(K^{1+\gamma})$ \\
Summation
  & $\sum_{j=1}^K j^{2\gamma}\sim K^{1+2\gamma}/(1+2\gamma)$
  & $C_{\mathrm{quad}}=O(K^{1+2\gamma})$ \\
Variational
  & \texttt{lem:kopt\_variational}
  & $K^*(c)=\frac{2\beta}{\alpha}\log c + \frac{2\gamma}{\alpha}\log\log c + O(1)$ \\
Layer-Survival
  & \texttt{thm:layer\_survival}
  & rate $O(c^{-\beta}(\log c)^{1+\gamma})$ \\
\bottomrule
\end{tabular}
\end{center}

\begin{remark}[Corrected leading coefficient for $K^*$]
\label{rem:kopt_fix}
Prior versions of this file contained the erroneous formula
\[
  K_{\mathrm{opt}} = \frac{2\beta}{\alpha(1+\gamma)}\log c
  \qquad\text{(obsolete pre-variational heuristic; superseded by
  \texttt{lem:kopt\_variational})}
\]
where an extraneous factor $(1+\gamma)^{-1}$ appeared from
a heuristic order-of-magnitude saddle balance.
The correct leading coefficient,
established in \texttt{lem:kopt\_variational}(ii) via
strict convexity + logarithmic FOC, is:
\[
  K^*(c) = \frac{2\beta}{\alpha}\log c + \frac{2\gamma}{\alpha}\log\log c + O(1).
\]
The Layer-Survival structure (dominance of $n=1$) does not depend
on the precise value of the leading coefficient,
but only on $K^*\asymp\log c$, which remains valid.
\end{remark}

Every step is now explicit.
The two remaining open items (O1--O4 in paper~XX) are:
\begin{enumerate}[label=(\roman*),nosep]
  \item Full Olver-Watson expansion for $Q_j$
    (near-turning-point integral, Section~\ref{sec:wkb}); and
  \item H2: $C_\infty < \infty$ for non-Airy models.
\end{enumerate}
Both are non-blocking for the conditional theorem.

\begin{thebibliography}{99}
\bibitem{Olver1997} F.~Olver,
  \textit{Asymptotics and Special Functions}, AK Peters, 1997.
\bibitem{Langer1949} R.~Langer,
  The asymptotic solutions of ordinary linear differential equations
  of the second order, with special reference to a turning point,
  \textit{Trans.\ Amer.\ Math.\ Soc.}\ \textbf{67} (1949), 461--490.
\bibitem{Olver1954} F.~Olver,
  The asymptotic solution of linear differential equations of the
  second order for large values of a parameter,
  \textit{Phil.\ Trans.\ Roy.\ Soc.\ London}\ \textbf{247} (1954),
  307--327.
\bibitem{BGT1987} N.~H.\ Bingham, C.~M.\ Goldie, J.~L.\ Teugels,
  \textit{Regular Variation},
  Cambridge University Press, 1987.
\end{thebibliography}

\end{document}
